\documentclass[11pt,a4paper]{article}

\usepackage{amsmath,amssymb,amsthm}
\usepackage{bm}
\usepackage{mathtools}

% Dirac notation and trace (replaces the physics package)
\newcommand{\Tr}{\operatorname{Tr}}
\newcommand{\ket}[1]{\left|#1\right\rangle}
\newcommand{\bra}[1]{\left\langle#1\right|}
\newcommand{\braket}[2]{\left\langle#1\middle|#2\right\rangle}
\usepackage{graphicx}
\usepackage{booktabs}
\usepackage{hyperref}
\usepackage{geometry}
\usepackage{microtype}
\usepackage{xcolor}

\geometry{margin=2.5cm}

\hypersetup{
  colorlinks=true,
  linkcolor=blue!70!black,
  citecolor=blue!70!black,
  urlcolor=blue!70!black,
}

% ---- convenience macros ----
\newcommand{\CC}{\mathbb{C}}
\newcommand{\RR}{\mathbb{R}}
\newcommand{\Hilbert}{\mathcal{H}}
\newcommand{\Kx}[2]{K_{#1,#2}}     % K_{x,w}
\newcommand{\Ex}[1]{E_{#1}}         % channel E_x
\newcommand{\Mx}[1]{M_{#1}}         % POVM element M_x
\newcommand{\NLL}{\mathcal{L}}

\theoremstyle{plain}
\newtheorem{proposition}{Proposition}

% ---- document ----
\title{\textbf{Quantum Channel Language Model (QCLM):\\
Architecture, Mathematics, and Implementation}}
\author{QCLM Project}
\date{\today}

\begin{document}

\maketitle
\tableofcontents
\bigskip

% ===========================================================================
\section{Overview}
% ===========================================================================

The \textbf{Quantum Channel Language Model (QCLM)} is an autoregressive
generative model whose sequential context-tracking mechanism is neither
self-attention (as in Transformers) nor a gated nonlinear recurrence (as in
RNNs, LSTMs, or GRUs), but a \emph{quantum channel} acting on a density
matrix in a finite-dimensional Hilbert space.

At each time step the model holds its entire context as a single
$n\times n$ complex density matrix $\rho_t$. Reading a token $x_t$ is
a \emph{measurement}: the next-token probability is given by the Born
rule $p(x_t \mid x_{<t}) = \Tr(\rho_{t-1}\, M_{x_t})$. Observing the
token \emph{collapses and evolves} the state through a
completely-positive, trace-preserving (CPTP) map --- a quantum channel.
There is no attention, no gate, and no nonlinear activation. The only
nonlinearity in the entire model is quantum measurement itself.

\paragraph{One-line identity.}
The QCLM is a \emph{complex-valued, trace-preserving linear state-space
model with a Born-rule (quadratic) readout} --- the quantum cousin of
modern linear SSMs (S4, Mamba), distinguished by genuine amplitude
\textbf{interference}.

% ===========================================================================
\section{Mathematical Formulation}
% ===========================================================================

\subsection{State space: the density matrix}

Let $\Hilbert = \CC^n$ be a Hilbert space of dimension $n$ (the
\emph{Hilbert-space dimension}, a hyperparameter). The model carries a
\emph{quantum state}
\[
  \rho \;\in\; \CC^{n\times n}, \qquad
  \rho = \rho^\dagger, \quad
  \rho \succeq 0, \quad
  \Tr(\rho) = 1,
\]
called a \emph{density matrix} (or density operator). A
\emph{pure state} takes the form $\rho = \ket{\psi}\!\bra{\psi}$ for
some unit vector $\ket{\psi}\in\CC^n$; a \emph{mixed state} is a convex
combination of pure states and encodes classical uncertainty over
quantum superpositions. The model begins each sequence in a learned
pure state $\rho_0 = \ket{\psi_0}\!\bra{\psi_0}$.

\subsection{Token operators: Kraus maps and POVM elements}

Each vocabulary token $x \in \{1,\ldots,V\}$ owns $W$ complex
$n\times n$ matrices called \textbf{Kraus operators}:
\[
  \bigl\{\Kx{x}{w}\bigr\}_{w=1}^{W}, \qquad \Kx{x}{w}\in\CC^{n\times n}.
\]
They define two objects simultaneously:

\begin{enumerate}
\item \textbf{Quantum channel} (completely-positive map):
\[
  \Ex{x}(\rho) \;=\; \sum_{w=1}^{W} \Kx{x}{w}\,\rho\,\Kx{x}{w}^\dagger.
\]

\item \textbf{POVM element} (positive operator):
\[
  \Mx{x} \;=\; \sum_{w=1}^{W} \Kx{x}{w}^\dagger \Kx{x}{w} \;\succeq\; 0.
\]
\end{enumerate}

\subsection{The isometry constraint}

A single global constraint ties all tokens together:
\begin{equation}\label{eq:isometry}
  \sum_{x=1}^{V}\sum_{w=1}^{W} \Kx{x}{w}^\dagger \Kx{x}{w} \;=\; I_n.
\end{equation}
Equation~\eqref{eq:isometry} is equivalent to $\sum_x \Mx{x} = I_n$,
which makes $\{\Mx{x}\}_{x=1}^{V}$ a valid \emph{positive
operator-valued measure} (POVM). This single equation is what
guarantees that next-token probabilities are exactly normalized --- by
construction, with no softmax.

\begin{proposition}[Exact normalization]
Under constraint~\eqref{eq:isometry}, for any density matrix $\rho$,
\[
  \sum_{x=1}^{V} \Tr(\rho\,\Mx{x})
  = \Tr\!\left(\rho \sum_x \Mx{x}\right)
  = \Tr(\rho\, I_n)
  = \Tr(\rho)
  = 1.
\]
\end{proposition}

\subsection{The autoregressive law (Born rule + collapse)}

Generation and likelihood scoring follow the standard postulates of
quantum measurement:

\begin{align}
  \text{Predict:} \quad &
    p(x_t \mid x_{<t}) \;=\; \Tr\!\bigl(\rho_{t-1}\,\Mx{x_t}\bigr)
    \label{eq:born}
  \\[6pt]
  \text{Update:} \quad &
    \rho_t \;=\; \frac{\Ex{x_t}(\rho_{t-1})}{p(x_t \mid x_{<t})}
    \label{eq:collapse}
\end{align}

Equation~\eqref{eq:born} is the \textbf{Born rule}: the probability of
outcome $x_t$ equals the expectation of the corresponding POVM element.
Equation~\eqref{eq:collapse} is the \textbf{post-measurement state
update}: apply the quantum channel $\Ex{x_t}$ and renormalize. Note
that
\[
  \Tr\bigl(\Ex{x_t}(\rho_{t-1})\bigr)
  = \Tr\!\bigl(\rho_{t-1}\,\Mx{x_t}\bigr)
  = p(x_t \mid x_{<t}),
\]
so the probability is just the trace of the (unnormalized) post-update
state --- it falls out for free, with no extra computation.

\subsection{Joint probability and the quantum forward algorithm}

Unrolling the recurrence, the joint probability of a string
$x_1,\ldots,x_T$ is
\begin{equation}\label{eq:joint}
  p(x_1,\ldots,x_T)
  = \Tr\!\Bigl(\Ex{x_T} \circ \cdots \circ \Ex{x_1}(\rho_0)\Bigr)
  = \Tr(\tilde\rho_T),
\end{equation}
where $\tilde\rho_T = \Ex{x_T}(\cdots(\Ex{x_1}(\rho_0))\cdots)$ is
the \emph{unnormalized} final state. Equivalently, the total
negative log-likelihood telescopes:
\begin{equation}\label{eq:nll_telescope}
  \NLL
  = -\sum_{t=1}^T \log p(x_t \mid x_{<t})
  = -\log \Tr(\tilde\rho_T),
\end{equation}
because each per-step log-probability is $\log\Tr(\tilde\rho_t) -
\log\Tr(\tilde\rho_{t-1})$, and the sum telescopes to
$\log\Tr(\rho_0) - \log\Tr(\tilde\rho_T) = -\log\Tr(\tilde\rho_T)$.
This is trained directly by gradient descent (the \emph{quantum forward
algorithm}).

\subsection{The role of quantum coherence}

Between measurements the density matrix carries \emph{off-diagonal}
entries $\rho_{ij}$ with $i\neq j$. These encode complex relative
phases between basis states. When the Born rule computes the next
probability,
\[
  p(x) = \Tr(\rho\,\Mx{x})
       = \sum_{i,j} \rho_{ij}\,(\Mx{x})_{ji},
\]
these off-diagonal terms contribute through \emph{quantum interference}:
different ``paths'' (Kraus operator sequences) can add or subtract in
amplitude before squaring. A classical model has $\rho_{ij}=0$ for
$i\neq j$ by construction, so it cannot exploit this interference.

% ===========================================================================
\section{Parametrization and Constraint Enforcement}
% ===========================================================================

\subsection{Free parameters}

The Kraus operators $\{\Kx{x}{w}\}$ are the model's trainable
parameters. We store them as separate real and imaginary parts:
\[
  \mathtt{Kr} \in \RR^{V\times W\times n\times n}, \qquad
  \mathtt{Ki} \in \RR^{V\times W\times n\times n},
\]
so that $\Kx{x}{w} = \mathtt{Kr}[x,w] + i\cdot\mathtt{Ki}[x,w]$. The
initial pure state is parametrized as a free complex vector
$\ket{\psi_0} = \mathtt{psi0\_r} + i\cdot\mathtt{psi0\_i}$,
normalized on each forward pass.

Total trainable parameters: $2\,V\,W\,n^2 + 2n$.

\subsection{Differentiable polar projection (isometry enforcement)}

The raw parameters do not in general satisfy the isometry
constraint~\eqref{eq:isometry}. We enforce it \emph{on every forward
pass} via a differentiable polar projection.

Stack all $V\cdot W$ Kraus operators row-wise:
\[
  \mathbf{V} \;=\;
  \begin{bmatrix}
    \Kx{1}{1} \\ \Kx{1}{2} \\ \vdots \\ \Kx{V}{W}
  \end{bmatrix}
  \;\in\; \CC^{(VWn)\times n}.
\]
The Gram matrix $G = \mathbf{V}^\dagger \mathbf{V} \in \CC^{n\times n}$
satisfies $G = \sum_{x,w}\Kx{x}{w}^\dagger\Kx{x}{w}$. We project:
\begin{equation}\label{eq:projection}
  \mathbf{V}_{\mathrm{iso}} \;=\; \mathbf{V}\,G^{-1/2},
\end{equation}
where $G^{-1/2}$ is computed via Hermitian eigendecomposition
$G = U\Lambda U^\dagger \Rightarrow G^{-1/2} = U\Lambda^{-1/2}U^\dagger$,
clamping small eigenvalues for numerical stability. By construction,
$\mathbf{V}_{\mathrm{iso}}^\dagger \mathbf{V}_{\mathrm{iso}} = I_n$,
so constraint~\eqref{eq:isometry} holds exactly.

The map $\mathbf{V} \mapsto \mathbf{V}_{\mathrm{iso}}$ is differentiable
(via \texttt{torch.linalg.eigh}), so gradients flow to the unconstrained
parameters through the projection.

% ===========================================================================
\section{Training}
% ===========================================================================

\subsection{Loss: negative log-likelihood in nats and bits}

The training objective is the sequence \textbf{negative log-likelihood}
(NLL) averaged over tokens:
\begin{equation}\label{eq:loss}
  \mathcal{L} \;=\; -\frac{1}{T}\sum_{t=1}^{T}\log p(x_t \mid x_{<t})
  \;\text{ (nats/token)}.
\end{equation}
This is the \emph{quantum forward algorithm}: at each step, the
probability $p(x_t\mid x_{<t}) = \Tr(\Ex{x_t}(\rho_{t-1}))$ is read
from the trace of the channel output, and the state is updated by
normalizing that output.

\paragraph{Bits per token.}
The equivalent metric in bits is
\begin{equation}\label{eq:bits}
  \text{bits/token} \;=\; \frac{\mathcal{L}}{\ln 2}
  \;=\; \log_2 \mathrm{PPL},
\end{equation}
where $\mathrm{PPL} = e^{\mathcal{L}}$ is perplexity. Bits per
token is the information-theoretic cross-entropy: it measures how many
bits the model needs on average to encode each token under its predicted
distribution. The theoretical minimum is the true entropy of the data
(unknown in practice); lower is better. The relation is simply a change
of logarithm base and carries no additional information beyond the NLL.

\subsection{Efficiency: training cost is independent of vocabulary size}

A key property follows from the isometry constraint. The Born
probability $p(x_t\mid x_{<t}) = \Tr\bigl(\Ex{x_t}(\rho_{t-1})\bigr)$
is the trace of the channel output \emph{for the observed token only}.
Computing this requires
\[
  \tilde\rho_t = \sum_{w=1}^W \Kx{x_t}{w}\,\rho_{t-1}\,\Kx{x_t}{w}^\dagger,
\]
which involves $W$ matrix multiplications of size $n\times n$. Because
$\Tr(\tilde\rho_t)$ is the probability and the state update is just
$\rho_t = \tilde\rho_t / \Tr(\tilde\rho_t)$, there is \textbf{no
need to evaluate the full vocabulary sum} during training. The per-token
training cost is $O(W n^3)$, independent of $V$.

This stands in contrast to transformer language models, where the final
softmax over the vocabulary is $O(V d)$ per token. For a vocabulary of
$V = 16{,}384$ and Hilbert dimension $n = 128$, the QCLM training cost
is $O(4 \cdot 128^3) \approx 8\text{M}$ operations per token, regardless
of vocabulary size.

\subsection{Truncated backpropagation through time (TBPTT)}

The quantum state $\rho_t$ depends on all previous tokens
$x_1,\ldots,x_t$ through the sequential recurrence. For long sequences,
retaining the full computation graph is memory-intensive. We apply
\textbf{truncated BPTT}: every $K$ steps the quantum state is detached
from the computation graph,
\[
  \rho_{tK} \;\leftarrow\; \texttt{stop\_gradient}(\rho_{tK}),
\]
so gradients are propagated through at most $K$ steps at a time.
This bounds activation memory without changing the forward recurrence
(the state evolution is unchanged; only the gradient path is truncated).

% ===========================================================================
\section{The Decoherence Ablation}
% ===========================================================================

To isolate the contribution of quantum coherence, we define a
\textbf{decohered} (classical) twin: after every state update, we zero
the off-diagonal entries of $\rho$,
\begin{equation}\label{eq:decohere}
  \rho_t^{\mathrm{cl}} \;=\;
  \mathrm{diag}\!\bigl(\rho_t^{\mathrm{diag}}\bigr)
  \;\Big/\;
  \Tr\!\bigl(\mathrm{diag}(\rho_t^{\mathrm{diag}})\bigr),
\end{equation}
where $\rho_t^{\mathrm{diag}}$ is the vector of diagonal entries of
$\rho_t$.

This operation:
\begin{enumerate}
\item \textbf{Exactly} reduces the QCLM to a classical $n$-state
  \emph{Hidden Markov Model} (HMM). The diagonal $(\rho_{11},\ldots,\rho_{nn})$
  plays the role of a classical probability vector over $n$ hidden states.
\item \textbf{Preserves} all parameters ($V,W,n$, Kraus operators).
  The decohered and quantum models have identical parameter counts and
  sizes.
\item Provides a \textbf{clean, dimension-matched ablation}: any
  performance difference between quantum and decohered models must come
  from the off-diagonal coherences, i.e., from quantum interference ---
  not from parameter count or architectural differences.
\end{enumerate}

% ===========================================================================
\section{Computational Structure and Engineering}
% ===========================================================================

\subsection{Per-token state update}

The hot path in \texttt{forward()} performs, at each step $t$:
\begin{enumerate}
\item \textbf{Channel application}: $\tilde\rho_t = \sum_w \Kx{x_t}{w}\,\rho_{t-1}\,\Kx{x_t}{w}^\dagger$ \\
  Implemented as two batched matrix multiplications:
  $\tilde{K}_w = \Kx{x_t}{w}\rho_{t-1}$, then $E = \sum_w \tilde{K}_w \Kx{x_t}{w}^\dagger$.
\item \textbf{Trace}: $p_t = \Tr(\tilde\rho_t) = \sum_i (\tilde\rho_t)_{ii}$ \quad (Born probability).
\item \textbf{Normalize}: $\rho_t = \tilde\rho_t / p_t$.
\item \textbf{NLL accumulation}: $\mathcal{L} \mathrel{-}= \log p_t$.
\end{enumerate}

\subsection{Real 2$\times$2 block matmuls for tensor-core utilization}

Modern NVIDIA GPUs (Ampere, Hopper, Blackwell) have tensor cores that
accelerate real \texttt{bfloat16}/\texttt{fp16} matrix multiplications
but do not natively accelerate \texttt{complex64}. We implement each
complex $(n\times n)@(n\times n)$ product as a single real
$(2n\times 2n)@(2n\times n)$ product in \texttt{bfloat16}:
\begin{equation}\label{eq:block}
  \begin{bmatrix}
    C_r \\ C_i
  \end{bmatrix}
  =
  \underbrace{
  \begin{bmatrix}
    A_r & -A_i \\
    A_i &  A_r
  \end{bmatrix}
  }_{\text{$A$-block, }(2n\times 2n)}
  \underbrace{
  \begin{bmatrix}
    B_r \\ B_i
  \end{bmatrix}
  }_{\text{$B$-block, }(2n\times n)},
\end{equation}
where $A = A_r + iA_i$ and $B = B_r + iB_i$ are the complex factors.
This replaces four separate real multiplications with one larger call,
which gives a single kernel launch and a larger tile for the tensor
cores to saturate. The \texttt{fp32} master weights in $\mathtt{Kr}$,
$\mathtt{Ki}$ are unaffected; only the matmul compute is in
\texttt{bfloat16}. The isometry projection stays in \texttt{float32}
and \texttt{complex64} for numerical stability (called once per forward
pass, not per token).

\subsection{Compiled chunked scan}

\paragraph{The Python-loop bottleneck.}
The default forward pass runs a Python \texttt{for t in range(L)} loop
over all $L$ tokens. Each iteration launches several GPU kernels and
incurs Python interpreter overhead ($\sim$10\,\textmu s/iteration),
amounting to $\sim$5\,ms for $L=512$ --- significant for fast GPUs.

\paragraph{Chunked compiled scan.}
We replace the $L$-iteration Python loop with:
\begin{enumerate}
\item \textbf{Pre-gather}: compute
  $\mathtt{Kr\_seq} = \mathtt{Kr\_iso}[\mathtt{tokens}]$ of shape
  $(B\times L\times W\times n\times n)$ in a \emph{single} batched
  index operation, instead of $L$ per-token gather kernels.
\item \textbf{Outer loop} over $\lceil L/K \rceil$ TBPTT chunks (e.g., 4
  iterations for $L=512$, $K=128$). At each chunk boundary the state is
  detached.
\item \textbf{Compiled inner scan}: each chunk body runs through
  \texttt{torch.compile(fullgraph=True, dynamic=False)}, which (a)
  unrolls the inner \texttt{for t in range(K)} loop at trace time,
  (b) fuses the block matmuls, trace, log, and divide across all $K$
  iterations into a small set of Triton kernels, and (c) eliminates
  Python interpreter overhead within the chunk entirely.
\end{enumerate}
The Python interpreter overhead drops from $\sim L$ to $\sim L/K$
iterations; the GPU sees a fused kernel stream rather than $L$ separate
launches.

\paragraph{Why not a parallel Blelloch scan?}
The composition $\Ex{x_t}\circ\cdots\circ\Ex{x_1}$ is associative, so
a parallel prefix scan exists in principle. However, its elements are
Liouville transfer matrices of size $n^2\times n^2$: for $n=128$, each
matrix is $16{,}384\times 16{,}384$ complex $\approx 4\,\text{GB}$.
This is memory-infeasible for any practically useful $n$. The chunked
compiled scan is the practical optimum.

% ===========================================================================
\section{Architecture Summary}
% ===========================================================================

Table~\ref{tab:arch} collects the key quantities.

\begin{table}[h]
\centering
\caption{QCLM architecture at a glance.}
\label{tab:arch}
\begin{tabular}{@{}lll@{}}
\toprule
\textbf{Quantity} & \textbf{Expression} & \textbf{Notes} \\
\midrule
State              & $\rho\in\CC^{n\times n}$ & density matrix, Hermitian, PSD, $\Tr\rho=1$ \\
Token operator     & $\{\Kx{x}{w}\}_{w=1}^W\subset\CC^{n\times n}$ & $W$ Kraus ops per token \\
Channel            & $\Ex{x}(\rho)=\sum_w\Kx{x}{w}\rho\Kx{x}{w}^\dagger$ & CPTP map \\
POVM element       & $\Mx{x}=\sum_w\Kx{x}{w}^\dagger\Kx{x}{w}$ & PSD \\
Constraint         & $\sum_{x,w}\Kx{x}{w}^\dagger\Kx{x}{w}=I_n$ & normalization \\
Prediction         & $p(x_t\mid x_{<t})=\Tr(\rho_{t-1}\Mx{x_t})$ & Born rule \\
State update       & $\rho_t = \Ex{x_t}(\rho_{t-1})/p(x_t\mid x_{<t})$ & collapse + renorm \\
Loss               & $\mathcal{L}=-\frac{1}{T}\sum_t\log\Tr(\Ex{x_t}(\rho_{t-1}))$ & nats/token \\
Bits/token         & $\mathcal{L}/\ln 2$ & $=\log_2\mathrm{PPL}$ \\
Parameters         & $2VWn^2 + 2n$ & real numbers; complex $= 2$ reals \\
Per-token cost     & $O(Wn^3)$ & \emph{independent of} $V$ \\
\bottomrule
\end{tabular}
\end{table}

% ===========================================================================
\section{Scaling}
% ===========================================================================

\subsection{Parameter count}

The dominant cost is the Kraus parameter table:
\[
  \text{params} \;\approx\; 2\,V\,W\,n^2,
\]
where the factor 2 counts real and imaginary parts. Parameters live
in the per-token Kraus operators, analogously to a joint embedding +
unembedding table. For the 2B target configuration:

\begin{center}
\begin{tabular}{@{}lrrrl@{}}
\toprule
\textbf{Config} & $V$ & $n$ & $W$ & \textbf{Params} \\
\midrule
PoC (char)         & 65    & 48  & 4 & $\approx 1.2\text{M}$ \\
1B (BPE-8k)        & 8{,}192  & 128 & 4 & $\approx 1.07\text{B}$ \\
2B (BPE-16k)       & 16{,}384 & 128 & 4 & $\approx 2.15\text{B}$ \\
\bottomrule
\end{tabular}
\end{center}

\subsection{Vocabulary is computationally cheap}

Because training cost is $O(Wn^3)$ per token (independent of $V$),
increasing the vocabulary from characters ($V=65$) to BPE subwords
($V\sim 10^4$) multiplies the parameter count by $V_{\text{new}}/V_{\text{old}}$
but does \emph{not} increase training FLOPs. The practical binding
constraints are memory bandwidth (loading the Kraus table) and the
$O(n^3)$ dense linear algebra per step.

\subsection{Scaling the Hilbert-space dimension}

Increasing $n$ increases both model capacity ($\propto n^2$ in params)
and per-token compute ($\propto n^3$). PoC experiments confirm that
the quantum advantage over the decohered HMM baseline \emph{grows}
with $n$:

\begin{center}
\begin{tabular}{@{}lrrr@{}}
\toprule
$n$ & Quantum (bits/char) & Classical (bits/char) & Advantage $\Delta$ \\
\midrule
24 & 3.008 & 3.208 & 0.200 \\
48 & 2.806 & 3.045 & 0.239 \\
\bottomrule
\end{tabular}
\end{center}

% ===========================================================================
\section{Connection to Prior Work}
% ===========================================================================

\paragraph{Hidden Quantum Markov Models (HQMMs).}
The QCLM is a generative HQMM \cite{monras2010,srinivasan2018}. Prior
work trained HQMMs on low-order toy sequences using EM. The QCLM differs
in (i) end-to-end gradient descent with a clean differentiable isometry
parametrization, (ii) application to natural language at scale, and
(iii) the controlled decoherence ablation that isolates interference.

\paragraph{Linear state-space models.}
S4 \cite{gu2021}, Mamba \cite{gu2023}, and related models use a linear
recurrence $h_t = A h_{t-1} + B u_t$. The QCLM's recurrence
$\rho_t = \Ex{x_t}(\rho_{t-1})/p(x_t)$ is also linear in $\rho_{t-1}$
(for fixed $x_t$), but the state is a \emph{density matrix} (complex,
Hermitian, PSD) and the readout is the Born rule (quadratic in
amplitudes) rather than a linear projection. The key distinction is the
capacity for \emph{genuine quantum interference} through off-diagonal
coherences.

\paragraph{Quantum machine learning.}
Parameterized quantum circuits (PQCs) are typically studied in the
context of near-term quantum hardware \cite{cerezo2021}. The QCLM is
a classical simulation of a quantum process, using quantum \emph{mathematics}
(not hardware) as the computational substrate, running efficiently on
CPUs and GPUs.

% ---- Bibliography (manual, no .bib file needed) ---------------------------
\begin{thebibliography}{9}

\bibitem{monras2010}
A.~Monras, A.~Beige, and K.~Wiesner.
\newblock Hidden quantum Markov models and non-adaptive read-out of many-body
  states.
\newblock \emph{Applied Mathematics \& Computational Science}, 1:93--107, 2010.

\bibitem{srinivasan2018}
S.~Srinivasan, G.~Gordon, and B.~Boots.
\newblock Learning hidden quantum Markov models.
\newblock In \emph{Proceedings of AISTATS}, 2018.

\bibitem{gu2021}
A.~Gu, K.~Goel, and C.~Ré.
\newblock Efficiently modeling long sequences with structured state spaces.
\newblock In \emph{Proceedings of ICLR}, 2022.

\bibitem{gu2023}
A.~Gu and T.~Dao.
\newblock Mamba: Linear-time sequence modeling with selective state spaces.
\newblock \emph{arXiv:2312.00752}, 2023.

\bibitem{cerezo2021}
M.~Cerezo et al.
\newblock Variational quantum algorithms.
\newblock \emph{Nature Reviews Physics}, 3:625--644, 2021.

\end{thebibliography}

\end{document}
